% =====================================================
% 题型：样本数据中位数选择题 (q_stats_median.tex)
% 知识板块：统计
% 主思维：特殊值验证方法
% 副思维：结构归纳
% 错因标签：未排序直接求中位数
% 讲义用途：基础检测题
% =====================================================
% 难度：★ (基础)
% =====================================================
\newcommand{\qStatsMedianStem}{样本数据 \(6, 8, 4, 5, 12\) 的中位数为 \hfill （\quad）}

\newcommand{\qStatsMedianOptions}{%
  \mcoptions[4]{\(5\)}{\(6\)}{\(8\)}{\(9\)}%
}

\newcommand{\qStatsMedianAnswer}{B}

\newcommand{\qStatsMedianSolution}{%
  【解题思路】\\
  求样本数据的中位数，第一步必须将原数据进行大小排序。排序后，若样本容量为奇数，则中位数为最中间位置的数据；若样本容量为偶数，则中位数为最中间两个数据的平均数。\\
  【解析计算】\\
  将原数据 \(6, 8, 4, 5, 12\) 从小到大排序为：\(4, 5, 6, 8, 12\)。\\
  因为数据个数为 \(5\)（奇数），所以中位数为第 \(3\) 个数，即 \(6\)。\\
  故选 B。%
}

\newcommand{\qStatsMedianDiagnosis}{%
  学生容易直接将未排序的原始数据中间项（即 4）误认为中位数。强调“排序”是求中位数的第一步。%
}

\newcommand{\qStatsMedianTeachingNotes}{%
  本题为基础题，重点强调排序概念。可引申至偶数个项时中位数的计算规则，以作防错点拨。%
}
